\chapter{测试与分析}
\section{测试环境与方法}
测试环境为 Windows、Python \SamplePythonVersion{} 和 SQLite \SampleSQLiteVersion{}。数据库初始化两个用户和三个座位，测试日期固定为 2026 年 10 月 9 日，跨日期用例使用次日。所有测试使用程序断言核对返回值，运行记录保存为 JSON 文件，并在论文编译前重新生成。

测试按固定顺序执行。首先建立基础预约，再检查重叠与相邻时段；随后验证输入、资源和取消归属；最后检验取消后的再次预约、跨日期预约和统计结果。这样可以同时检查单次操作及操作之间的状态衔接。

\section{功能测试结果}
实际执行的用例与结果见表\ref{tab:test-cases}。起止时刻均以同一日期内的本地时间表示。

\begingroup
\tablefont
\begin{longtable}{@{}p{13mm}p{71mm}p{43mm}p{13mm}@{}}
  \caption{功能测试用例与执行结果}\label{tab:test-cases}\\
  \toprule
  编号 & 操作或条件 & 预期结果 & 结果\\
  \midrule
  \endfirsthead
  \multicolumn{4}{r}{\heiti 续表\thetable}\\
  \toprule
  编号 & 操作或条件 & 预期结果 & 结果\\
  \midrule
  \endhead
  \bottomrule
  \endfoot
  T01 & 用户 1 预约座位 1，09:00 至 10:00 & 创建成功 & 通过\\
  T02 & 用户 2 预约同一座位，09:30 至 10:30 & 座位冲突 & 通过\\
  T03 & 用户 2 预约同一座位，10:00 至 11:00 & 创建成功 & 通过\\
  T04 & 用户 1 预约座位 2，09:30 至 10:30 & 用户时段冲突 & 通过\\
  T05 & 提交结束早于开始的时段 & 输入无效 & 通过\\
  T06 & 提交不存在的座位编号 & 座位不存在 & 通过\\
  T07 & 用户 2 取消用户 1 的预约 & 取消失败 & 通过\\
  T08 & 用户 1 取消本人有效预约 & 取消成功 & 通过\\
  T09 & 用户 2 预约已释放的 09:00 至 10:00 & 创建成功 & 通过\\
  T10 & 再次取消同一条已取消记录 & 取消失败 & 通过\\
  T11 & 用户 1 预约次日同座位与时段 & 创建成功 & 通过\\
  T12 & 查询首日有效预约数量与分钟数 & 2 条，共 120 分钟 & 通过\\
\end{longtable}
\endgroup

\section{结果分析}
\subsection{业务规则验证}
\SampleTestCount{} 项用例中，\SamplePassedCount{} 项通过。T02 与 T03 验证了重叠时段受到限制而相邻时段可以共存；T04 验证了同一用户的时间限制；T07、T08 与 T10 验证了取消归属和重复操作；T09 说明取消记录不再占用资源；T11 说明冲突检查正确区分日期。测试运行结束后，首日仍有两条有效预约，累计预约时长为 120 分钟。

\subsection{预约时长统计}
为演示统计口径，以预约时长占全部可预约时长的比例定义预约占用率 $U$，计算方法见式\eqref{eq:utilization}。
\begin{equation}
  U = \frac{\displaystyle\sum_{i=1}^{n}(e_i-s_i)}{S T}\times 100\%
  \label{eq:utilization}
\end{equation}
其中，$n$ 为当日有效预约数，$S$ 为开放座位数，$T$ 为单个座位每日可预约分钟数。本组数据中，$S=3$、$T=840$，有效预约共 120 分钟，故 $U=\SampleUtilization{}\%$。该指标反映记录中的预约时长；实际到场使用情况需要由签到数据进一步确认。

\section{测试范围与后续验证}
本次用例验证了小规模串行业务的主要路径。后续应补充开放时段端点、无效日期、未知用户、数据库重启与多连接同时预约等情形，并检查锁等待、失败重试和日志记录。采用文件数据库和实际用户入口后，还需要开展持久化与身份认证测试。
